\documentclass[a4paper,landscape]{article}
\usepackage{puzzlebug}
\begin{document}
\begin{center}
\begin{tikzpicture}[line cap=round,line join=round]
\clip (0,0) rectangle (29.7,21);
\fill[white] (0,0) rectangle (29.7,21);
\foreach \x in {7.425,14.85,22.275}
    \draw[red] (\x,0) -- +(0, 21);
\draw[red] (0,10.5) -- +(29.7,0);


%%%%%%%%%%%%%%%%%%%% FRONT COVER %%%%%%%%%%%%%%%%%%%%
\begin{scope}[shift={(7.425,0)}]

\clip (0,0) rectangle (7.425,10.5);

\draw[line width=1pt,fill=white]
    (3.1,7.4) -- (3.7,8.4) [rounded corners=5pt] -- (6.6,8.4) -- (6.6,9.6)
    -- (0.824,9.6) -- (0.824,8.4) [rounded corners=0pt]
    -- (3.5,8.4) -- cycle;
\node at (3.7125,9) {\setmainfont[Scale=2.5]{Bubble Sans}Puzzlebug};
\node[anchor=north east] at (6.6,8.4) {\setmainfont[Scale=1.3]{Bubble Sans}\#2};
\node[anchor=east,rotate=90] at (6.8,9.7) {\setmainfont[Scale=0.5]{Bubble Sans}created by Matthew Scroggs};


\begin{scope}[shift={(1,7)}]
\puzzlebug[large]{r}{(0,0)}{(0.02,0)}{neutral}
\end{scope}

\node[anchor=south,text width=7cm] at (3.7125,0.5) {
\null\hfill\scalebox{0.8}{
\parbox{7.5cm}{\centering
{\setmainfont[Scale=1.8]{Bangers}Featuring}

\par\smallskip


\mbox{names}
\quad
\mbox{of}
\quad
\mbox{puzzles}
\quad
\mbox{go}
\quad
\mbox{here}

}}\hfill\null};


\end{scope}

%%%%%%%%%%%%%%%%%%%% PAGE ONE %%%%%%%%%%%%%%%%%%%%
\begin{scope}[shift={(14.85,0)}]

\puzztitle{evil counts}

\begin{scope}[shift={(0.9625,1.5)},scale=0.61]
\tiny
\fill[white] (-1,0) rectangle +(11,11);
\fill[black] (-1,0) rectangle +(1,1) (-1,1) rectangle +(1,1) (-1,2) rectangle +(1,1) (-1,3) rectangle +(1,1) (-1,4) rectangle +(1,1) (-1,5) rectangle +(1,1) (-1,6) rectangle +(1,1) (-1,7) rectangle +(1,1) (-1,8) rectangle +(1,1) (-1,9) rectangle +(1,1) (-1,10) rectangle +(1,1)(0,10) rectangle +(1,1) (1,10) rectangle +(1,1) (2,10) rectangle +(1,1) (3,10) rectangle +(1,1) (4,10) rectangle +(1,1) (5,10) rectangle +(1,1) (6,10) rectangle +(1,1) (7,10) rectangle +(1,1) (8,10) rectangle +(1,1) (9,10) rectangle +(1,1)(0,9) rectangle +(1,1) (1,9) rectangle +(1,1) (9,9) rectangle +(1,1) (3,8) rectangle +(1,1) (9,8) rectangle +(1,1) (3,7) rectangle +(1,1) (6,7) rectangle +(1,1) (2,6) rectangle +(1,1) (3,6) rectangle +(1,1) (4,6) rectangle +(1,1) (6,6) rectangle +(1,1) (7,6) rectangle +(1,1) (8,6) rectangle +(1,1) (6,5) rectangle +(1,1) (3,4) rectangle +(1,1) (1,3) rectangle +(1,1) (2,3) rectangle +(1,1) (3,3) rectangle +(1,1) (5,3) rectangle +(1,1) (6,3) rectangle +(1,1) (7,3) rectangle +(1,1) (3,2) rectangle +(1,1) (6,2) rectangle +(1,1) (0,1) rectangle +(1,1) (6,1) rectangle +(1,1) (0,0) rectangle +(1,1) (8,0) rectangle +(1,1) (9,0) rectangle +(1,1);
\fill[black!20!white] (2,9) -- +(0,1) -- +(-1,1) -- cycle;
\node[inner sep=2pt,anchor=east] at (2,9.7) {OC};
\fill[black!20!white] (0,8) -- +(0,1) -- +(-1,1) -- cycle;
\node[inner sep=2pt,anchor=east] at (0,8.7) {SO};
\fill[black!20!white] (4,8) -- +(0,1) -- +(-1,1) -- cycle;
\node[inner sep=2pt,anchor=east] at (4,8.7) {SS};
\fill[black!20!white] (0,7) -- +(0,1) -- +(-1,1) -- cycle;
\node[inner sep=2pt,anchor=east] at (0,7.7) {L};
\fill[black!20!white] (4,7) -- +(0,1) -- +(-1,1) -- cycle;
\node[inner sep=2pt,anchor=east] at (4,7.7) {O};
\fill[black!20!white] (7,7) -- +(0,1) -- +(-1,1) -- cycle;
\node[inner sep=2pt,anchor=east] at (7,7.7) {T};
\fill[black!20!white] (0,6) -- +(0,1) -- +(-1,1) -- cycle;
\node[inner sep=2pt,anchor=east] at (0,6.7) {VU};
\fill[black!20!white] (0,5) -- +(0,1) -- +(-1,1) -- cycle;
\node[inner sep=2pt,anchor=east] at (0,5.7) {UI};
\fill[black!20!white] (7,5) -- +(0,1) -- +(-1,1) -- cycle;
\node[inner sep=2pt,anchor=east] at (7,5.7) {VS};
\fill[black!20!white] (0,4) -- +(0,1) -- +(-1,1) -- cycle;
\node[inner sep=2pt,anchor=east] at (0,4.7) {VS};
\fill[black!20!white] (4,4) -- +(0,1) -- +(-1,1) -- cycle;
\node[inner sep=2pt,anchor=east] at (4,4.7) {UL};
\fill[black!20!white] (8,3) -- +(0,1) -- +(-1,1) -- cycle;
\node[inner sep=2pt,anchor=east] at (8,3.7) {N};
\fill[black!20!white] (0,2) -- +(0,1) -- +(-1,1) -- cycle;
\node[inner sep=2pt,anchor=east] at (0,2.7) {SU};
\fill[black!20!white] (4,2) -- +(0,1) -- +(-1,1) -- cycle;
\node[inner sep=2pt,anchor=east] at (4,2.7) {VO};
\fill[black!20!white] (7,2) -- +(0,1) -- +(-1,1) -- cycle;
\node[inner sep=2pt,anchor=east] at (7,2.7) {L};
\fill[black!20!white] (1,1) -- +(0,1) -- +(-1,1) -- cycle;
\node[inner sep=2pt,anchor=east] at (1,1.7) {VE};
\fill[black!20!white] (7,1) -- +(0,1) -- +(-1,1) -- cycle;
\node[inner sep=2pt,anchor=east] at (7,1.7) {I};
\fill[black!20!white] (1,0) -- +(0,1) -- +(-1,1) -- cycle;
\node[inner sep=2pt,anchor=east] at (1,0.7) {ST};
\fill[black!20!white] (0,9) -- +(1,0) -- +(0,1) -- cycle;
\node[inner sep=2pt,anchor=south] at (0.3,9) {UC};
\fill[black!20!white] (1,9) -- +(1,0) -- +(0,1) -- cycle;
\node[inner sep=2pt,anchor=south] at (1.3,9) {SU};
\fill[black!20!white] (1,3) -- +(1,0) -- +(0,1) -- cycle;
\node[inner sep=2pt,anchor=south] at (1.3,3) {VV};
\fill[black!20!white] (2,10) -- +(1,0) -- +(0,1) -- cycle;
\node[inner sep=2pt,anchor=south] at (2.3,10) {VC};
\fill[black!20!white] (2,6) -- +(1,0) -- +(0,1) -- cycle;
\node[inner sep=2pt,anchor=south] at (2.3,6) {VL};
\fill[black!20!white] (2,3) -- +(1,0) -- +(0,1) -- cycle;
\node[inner sep=2pt,anchor=south] at (2.3,3) {VU};
\fill[black!20!white] (3,2) -- +(1,0) -- +(0,1) -- cycle;
\node[inner sep=2pt,anchor=south] at (3.3,2) {VV};
\fill[black!20!white] (4,10) -- +(1,0) -- +(0,1) -- cycle;
\node[inner sep=2pt,anchor=south] at (4.3,10) {I};
\fill[black!20!white] (4,6) -- +(1,0) -- +(0,1) -- cycle;
\node[inner sep=2pt,anchor=south] at (4.3,6) {SS};
\fill[black!20!white] (5,10) -- +(1,0) -- +(0,1) -- cycle;
\node[inner sep=2pt,anchor=south] at (5.3,10) {ST};
\fill[black!20!white] (5,3) -- +(1,0) -- +(0,1) -- cycle;
\node[inner sep=2pt,anchor=south] at (5.3,3) {SV};
\fill[black!20!white] (6,10) -- +(1,0) -- +(0,1) -- cycle;
\node[inner sep=2pt,anchor=south] at (6.3,10) {N};
\fill[black!20!white] (7,10) -- +(1,0) -- +(0,1) -- cycle;
\node[inner sep=2pt,anchor=south] at (7.3,10) {VN};
\fill[black!20!white] (7,6) -- +(1,0) -- +(0,1) -- cycle;
\node[inner sep=2pt,anchor=south] at (7.3,6) {VC};
\fill[black!20!white] (7,3) -- +(1,0) -- +(0,1) -- cycle;
\node[inner sep=2pt,anchor=south] at (7.3,3) {L};
\fill[black!20!white] (8,10) -- +(1,0) -- +(0,1) -- cycle;
\node[inner sep=2pt,anchor=south] at (8.3,10) {VT};
\fill[black!20!white] (8,6) -- +(1,0) -- +(0,1) -- cycle;
\node[inner sep=2pt,anchor=south] at (8.3,6) {VL};
\fill[black!20!white] (9,8) -- +(1,0) -- +(0,1) -- cycle;
\node[inner sep=2pt,anchor=south] at (9.3,8) {OV};
\foreach \x in {-1,...,10}
  \draw (\x,0) -- (\x,11);
\foreach \y in {0,...,11}
  \draw (-1,\y) -- (10,\y);
\draw (0,11) -- +(1,-1) (1,11) -- +(1,-1) (2,11) -- +(1,-1) (3,11) -- +(1,-1) (4,11) -- +(1,-1) (5,11) -- +(1,-1) (6,11) -- +(1,-1) (7,11) -- +(1,-1) (8,11) -- +(1,-1) (9,11) -- +(1,-1)(0,10) -- +(1,-1) (1,10) -- +(1,-1) (9,10) -- +(1,-1) (3,9) -- +(1,-1) (9,9) -- +(1,-1) (3,8) -- +(1,-1) (6,8) -- +(1,-1) (2,7) -- +(1,-1) (3,7) -- +(1,-1) (4,7) -- +(1,-1) (6,7) -- +(1,-1) (7,7) -- +(1,-1) (8,7) -- +(1,-1) (6,6) -- +(1,-1) (3,5) -- +(1,-1) (1,4) -- +(1,-1) (2,4) -- +(1,-1) (3,4) -- +(1,-1) (5,4) -- +(1,-1) (6,4) -- +(1,-1) (7,4) -- +(1,-1) (3,3) -- +(1,-1) (6,3) -- +(1,-1) (0,2) -- +(1,-1) (6,2) -- +(1,-1) (0,1) -- +(1,-1) (8,1) -- +(1,-1) (9,1) -- +(1,-1);
\node[black!50!white] at (5.5,6.5) {\large O};
\node[black!50!white] at (3.5,5.5) {\large L};
\node[black!50!white] at (6.5,4.5) {\large E};
\node[black!50!white] at (4.5,3.5) {\large U};

\node[black!50!white] at (2.5,9.5) {\color{red}\large V};
\node[black!50!white] at (3.5,9.5) {\color{red}\large I};
\node[black!50!white] at (4.5,9.5) {\color{red}\large O};
\node[black!50!white] at (5.5,9.5) {\color{red}\large L};
\node[black!50!white] at (6.5,9.5) {\color{red}\large E};
\node[black!50!white] at (7.5,9.5) {\color{red}\large N};
\node[black!50!white] at (8.5,9.5) {\color{red}\large T};
\end{scope}

\begin{scope}[shift={(5,9)}]
\puzzlebug{r}{(0,-0.025)}{(-0.01,-0.035)}{neutral}
\end{scope}

\pagenum{1}
\end{scope}

%%%%%%%%%%%%%%%%%%%% PAGE TWO %%%%%%%%%%%%%%%%%%%%
\begin{scope}[shift={(22.275,0)}]

\draw[line width=1pt,fill=white]
    (-0.5,9) -- (0.5,9.1) [rounded corners=5pt] -- (0.5,10.15)
    -- (6.9,10.15) -- (6.9,6.31) -- (0.5,6.31) [rounded corners=0pt]
    -- (0.5,8.8) -- cycle;


\node[anchor=north,text width=7cm] at (3.7125,10.08) {
\null\hfill\scalebox{0.6}{
\parbox{10.3cm}{\raggedright
\noindent
the letters e, v, i, l, c, o, u, n, t and s represent the digits 0 to 9, with each letter representing a different digit.
write a letter representing a non-zero digit in each white square so that the
sum of the digits in each horizontal and vertical run is the total given above or to the left of the run, with no letter repeated in the run.
for example, if s and o represent 1 and 2, then the three boxes in the second row next to the clue `SO' must all contain three different letters that represent digits adding to 12.
none of the numbers given as a clue starts with a leading 0.
when the puzzle is complete, the letters in the top row will spell out a word.
}}\hfill\null};


\puzztitle[5.2]{on both sides};

\begin{scope}[shift={(0,-4.8)}]
\begin{scope}[shift={(0,0.2)}]
\node[scale=0.6,anchor=south] at (0.8,8.7) {RAKE};
\node[scale=0.6,anchor=south] at (1.6,8.7) {THOS};
\node[scale=0.6,anchor=south] at (2.4,8.7) {ACH};
\node[scale=0.6,anchor=south] at (3.2,8.7) {OUTE};
\draw[line width=1pt,line join=round,line cap=round]
    (0.4,7.8) rectangle +(0.8*4,0.8)
    (1.2,7.8) -- +(0,0.8)
    (2.0,7.8) -- +(0,0.8)
    (2.8,7.8) -- +(0,0.8)
;

\node[black!50!white] at (0.8,8.2) {\color{red}\large D};
\node[black!50!white] at (1.6,8.2) {\color{red}\large E};
\node[black!50!white] at (2.4,8.2) {\color{red}\large E};
\node[black!50!white] at (3.2,8.2) {\color{red}\large R};

\end{scope}

\begin{scope}[shift={(2,6.75)}]
\puzzlebug{r}{(0,-0.02)}{(-0.01,0.025)}{neutral}
\end{scope}

\draw[line width=1pt,fill=white]
    (4,6.9) -- (4.35,7.18) [rounded corners=5pt] -- (4.1,7.18) -- (4.1,10.15)
    -- (7.1,10.15) -- (7.1,7.18) [rounded corners=0pt]
    -- (4.6,7.18) -- cycle;

\node[anchor=south,text width=3.5cm] at (5.62,7.2) {
\null\hfill\scalebox{0.6}{
\parbox{4.7cm}{\raggedright
\noindent
in each box, write the letter than can be put before and after the letters above the box
to make two words.
For example, you could write Q in the first box if QRAKE and RAKEQ were words.
The four letters you write will spell a word.
}}\hfill\null};

\end{scope}

\pagenum{2}
\end{scope}

%%%%%%%%%%%%%%%%%%%% PAGE THREE %%%%%%%%%%%%%%%%%%%%
\begin{scope}[shift={(29.7,21)},rotate=180,transform shape]


\pagenum{3}
\end{scope}

%%%%%%%%%%%%%%%%%%%% PAGE FOUR %%%%%%%%%%%%%%%%%%%%
\begin{scope}[shift={(22.275,21)},rotate=180]
\pagenum{4}
\end{scope}

%%%%%%%%%%%%%%%%%%%% PAGE FIVE %%%%%%%%%%%%%%%%%%%%
\begin{scope}[shift={(14.85,21)},rotate=180]
\pagenum{5}
\end{scope}

%%%%%%%%%%%%%%%%%%%% PAGE SIX %%%%%%%%%%%%%%%%%%%%
\begin{scope}[shift={(7.425,21)},rotate=180,transform shape]

\puzztitle{...and finally};

\begin{scope}[shift={(-0.3,0.4)}]
\node[anchor=south,scale=0.6,align=center] at (6.5,8.8) {extra\\letters};
\node[anchor=east,scale=0.6,align=right,red] at (2.25,8.5) {weight};
\node[anchor=east,scale=0.6,align=right,red] at (2.25,7.9) {alert};
\node[anchor=east,scale=0.6,align=right] at (2.25,7.3) {word from\\on both sides};
\node[anchor=east,scale=0.6,align=right,red] at (2.25,6.7) {lodge};
\node[anchor=east,scale=0.6,align=right] at (2.25,6.1) {top row of\\crypturo};
\node[red] at (2.6,6.1) {V};
\node[red] at (3.2,6.1) {I};
\node[red] at (3.8,6.1) {O};
\node[red] at (4.4,6.1) {L};
\node[red] at (5.0,6.1) {E};
\node[red] at (5.6,6.1) {T};
\node[red] at (6.5,6.1) {N};
\node[red] at (2.6,7.3) {R};
\node[red] at (3.2,7.3) {E};
\node[red] at (3.8,7.3) {D};
\node[red] at (6.5,7.3) {E};
\draw[line width=1pt,line join=round,line cap=round,shift={(2.3,5.8)},scale=0.6]
    (0,4) rectangle +(5,1)
    (1,4) -- +(0,1)
    (2,4) -- +(0,1)
    (3,4) -- +(0,1)
    (4,4) -- +(0,1)
    (0,3) rectangle +(4,1)
    (1,3) -- +(0,1)
    (2,3) -- +(0,1)
    (3,3) -- +(0,1)
    (0,2) rectangle +(3,1)
    (1,2) -- +(0,1)
    (2,2) -- +(0,1)
    (0,1) rectangle +(4,1)
    (1,1) -- +(0,1)
    (2,1) -- +(0,1)
    (3,1) -- +(0,1)
    (0,0) rectangle +(6,1)
    (1,0) -- +(0,1)
    (2,0) -- +(0,1)
    (3,0) -- +(0,1)
    (4,0) -- +(0,1)
    (5,0) -- +(0,1)
    (6.5,0) rectangle +(1,5)
    (6.5,1) -- +(1,0)
    (6.5,2) -- +(1,0)
    (6.5,3) -- +(1,0)
    (6.5,4) -- +(1,0)
;
\end{scope}

\draw[line width=1pt,fill=white]
    (5,1.2) -- (4.7,1.9) [rounded corners=5pt] -- (6.9,1.9) -- (6.9,4.1)
    -- (0.5,4.1) -- (0.5,1.4) -- (4,1.4) -- (4,1.9) [rounded corners=0pt]
    -- (4.5,1.9) -- cycle;

\node[anchor=north,text width=7cm] at (3.7125,4.08) {
\null\hfill\scalebox{0.6}{
\parbox{10.3cm}{\raggedright
\noindent The words that are
{\color{red}???},
{\color{red}???},
the answer to on both sides on page 2,
{\color{red}???} and
in the top row of the solution to evil counts on page 1
are all anagrams of themed words with an extra letter added. Write the themed words in the boxes above and the extra letters
in the boxes to the right. One you have done this, the extra letters will spell another word with the same theme:
you can use that word as a
discount code to get 50\% off Puzzlebug \#3.
}}\hfill\null};

\begin{scope}[shift={(6.7,1)}]
\puzzlebug{l}{(0,0)}{(0.02,0.02)}{neutral}
\end{scope}

\pagenum{6}
\end{scope}

%%%%%%%%%%%%%%%%%%%% BACK PAGE %%%%%%%%%%%%%%%%%%%%
\begin{scope}

\howtofold

\node[anchor=south,text width=7cm] at (3.7125,0.15) {
\null\hfill\scalebox{0.5}{
\parbox{12.5cm}{\raggedright
\begin{center}
\scalebox{1.3}{{\setmainfont[Scale=1.3]{Chalkdust Icon Font}d} mscroggs.co.uk/puzzlebug}
\end{center}

\vspace{3mm}

\noindent Puzzlebug \#2 was published in January 2026.
it was written and drawn by Matthew Scroggs.
Lettering is done using the fonts
Bubble Sans by Abay Emes
and
Bangers by Vernon Adams:
both are available under the SiL Open Font License version 1.1.

\begin{center}
\textcopyright{} Matthew Scroggs, 2026\\
iSSN 2978-4255 (print). iSSN 2978-4263 (online).\\
Published by mscroggs.co.uk, King's Lynn, UK.\\
\bugcode%
{1}% numbering version
{2}% issue number
{0}% edit number (0=draft, 1=first release)
{0}% 0=test, 1=print, 2=pdf
{0}% print run number
{0}% print run size
\end{center}
}}\hfill\null};


\end{scope}

\end{tikzpicture}
\end{center}
\end{document}
